\section{Lean formalization and object-language route}\label{sec:lean}

\subsection{Repository interfaces reused}

The ProveIt repository has two distinct levels that must not be conflated.
The ambient Lean metatheory provides convenient finite sums, rationals, and
proof automation.  The internal target is an explicit derivation
\[
 ([]:\operatorname{List}(\mathrm{Fm}\;\mathrm{PTm}\;k))
 \;\vdash[\mathrm{PRAAx}]\;\varphi
\]
in the formalized \(\ipra\) calculus.  The existing
\code{pra_iff_native} theorem then transports a closed \code{PRAAx} derivation
to the concise native scalar core.

The checked arithmetic library already contains the principal combinators
needed here:
\begin{itemize}
\item \code{IPRADerivations.Core.BoundedSums}: elimination of terms from a
  zero nonnegative sum and construction of a zero sum from pointwise zeros;
\item \code{IPRADerivations.Core.SumLemmas}: sum shift, sum of pointwise sums,
  bounded congruence, and ``a term is at most a nonnegative sum'';
\item \code{IPRADerivations.Core.MultiSums}: rectangular finite Fubini and
  reindexing;
\item \code{IPRADerivations.Core.Pairing}: coding product outcomes;
\item arithmetic, order, truncated subtraction, division with remainder, and
  ring-normalization modules used by \code{deriveNat};
\item \code{ConciseIPRA.NativeBiInterpretation}: the
  \code{pra_iff_native} bridge.
\end{itemize}

The earlier constructive central-limit packet supplies compatible design
patterns: a natural-weight \code{FinLaw}, nested rational interval
\code{Regular} names, and an \code{InternalProof} structure separating a
primitive-recursive checker from its actual \code{PRAAx} derivation.  The
present packet follows those patterns but does not assume the uncompleted CLT
blueprints as theorems.

\subsection{Ambient finite API}

The file \code{lean/ExpectedValueFinite.lean} begins with

\begin{lstlisting}[language={}]
structure FinLaw where
  card : Nat
  weight : Fin card -> Nat
  total : Nat
  total_pos : 0 < total
  weight_sum : (sum i, weight i) = total

def FinLaw.expect (L : FinLaw)
    (X : Fin L.card -> Rat) : Rat :=
  (sum i, (L.weight i : Rat) * X i) / L.total
\end{lstlisting}

It also defines event probability, moments, variance, finite CDFs, almost-sure
equality, pushforward fiber weights, product-law data, partition cell weights,
and conditional means.  The structures
\code{FiniteExpectationProofPackage},
\code{FiniteDependenceProofPackage}, and \code{SampleProofPackage} list the
ambient theorems to be constructed.  They are records of obligations, not
axioms: no inhabitant is provided.

A good ambient implementation should avoid a single enormous proof.  The
recommended module graph is

\begin{verbatim}
IPRADerivations/Probability/FiniteLaw.lean
IPRADerivations/Probability/Expectation.lean
IPRADerivations/Probability/TailSum.lean
IPRADerivations/Probability/Variance.lean
IPRADerivations/Probability/MarkovChebyshev.lean
IPRADerivations/Probability/CauchySchwarz.lean
IPRADerivations/Probability/Independence.lean
IPRADerivations/Probability/Conditional.lean
IPRADerivations/Probability/SampleMean.lean
IPRADerivations/Probability/Kolmogorov.lean
\end{verbatim}

Each file should first prove the theorem over ordinary \code{Rat} and finite
sums.  More importantly, it should state a denominator-cleared lemma whose
shape already matches the natural-number object language.

\subsection{Natural encodings}

Introduce codes such as

\begin{lstlisting}[language={}]
structure ZCode where
  pos neg : Nat

structure QCode where
  num : ZCode
  den : Nat
  denPos : 0 < den

structure FinLawCode where
  card total commonDen : Nat
  weightCode valueCode : Nat
\end{lstlisting}

A finite list is encoded by the repository's pairing/list scheme.  Access
outside the declared length returns zero.  The decoder validates:
\begin{enumerate}
\item positive total and positive common denominator;
\item bounded weight sum equals the total;
\item every observable entry has the declared common denominator;
\item theorem-specific side conditions, such as nonnegativity or a positive
  threshold.
\end{enumerate}
Invalid inputs return zero from the checker.  Thus a universally quantified
open equation \code{check(args)=0} acts as an implication: malformed data pass
vacuously, while valid data trigger the mathematical verification.

The essential primitive-recursive functions are
\begin{verbatim}
weightSum       weightedNum      eventMass
absWeightedNum secondMomentNum  fiberMass
integerTailMass productIndex    cellMass
firstCrossing  finiteMoment     regularIntervalCheck
\end{verbatim}
Their defining equations must be proved once and reused.  No semantic theorem
should unfold a large decoder repeatedly.

\subsection{Object-language theorem shapes}

The file \code{lean/ObjectLanguageTargets.lean} defines a generic open target

\begin{lstlisting}[language={}]
def checkedOpen {k : Nat} (check : PR k) : Fm PTm k :=
  PTm.app check Syn.var = PTm.zero

structure InternalProof {k : Nat}
    (semantic : (Fin k -> Nat) -> Nat)
    extends Representation semantic where
  derived : ([] : List (Fm PTm k))
    |- [PRAAx] checkedOpen check
\end{lstlisting}

The concrete syntax in the source uses ProveIt's equality and derivability
symbols; the ASCII display above is only schematic.  A completed theorem
contains both:
\begin{enumerate}
\item \code{correct}: the \code{PR} term evaluates to the intended semantic
  test; and
\item \code{derived}: an actual derivation proving that the test always
  returns zero.
\end{enumerate}
Standard-model truth establishes neither field automatically.

\subsection{Proof plan for each finite theorem}

\paragraph{Linearity.}
Normalize \(X,Y\) to a common denominator.  The numerator for \(X+Y\) is the
bounded sum of the pointwise numerator sum.  Instantiate \code{sum_add}, apply
the ring tool to denominator products, and close universally.

\paragraph{Monotonicity.}
The checker stores nonnegative numerator differences
\(d_i=Y_i-X_i\).  A bounded sum of nonnegative naturals is nonnegative.
Equivalently, define the target as the existence of a natural slack
\(s\) with \(\operatorname{num}(Y)-\operatorname{num}(X)=s\); the sum of the
pointwise slacks supplies \(s\).

\paragraph{Indicator and finite LOTUS.}
Indicator expectation unfolds to \code{eventMass}.  For LOTUS, define
\code{fiberMass(v)} as a bounded sum.  Apply \code{MultiSums} to the matrix
\(w_i\ind_{f(i)=v}g(v)\), and use the proof that exactly one finite fiber
indicator is active for each \(i\).

\paragraph{Tail sum.}
First derive the open identity
\[
 x=\sum_{k<B}\ind_{k<x}\qquad(x\le B)
\]
by quantifier-free induction on \(x\) or \(B\).  Substitute each encoded value,
multiply by the weight, and use finite Fubini.

\paragraph{Variance.}
After common-denominator normalization, invoke ring normalization on the
identity
\[
 W\sum_iw_i(X_i-\mu)^2
 =W\sum_iw_iX_i^2-(\sum_iw_iX_i)^2.
\]
The displayed form avoids division and is polynomial over signed-natural
codes.

\paragraph{Markov and Chebyshev.}
Compile the pointwise slack in
\(a\ind_{a\le X_i}\le X_i\), sum it, and clear denominators.  Chebyshev is the
same theorem applied to the encoded square deviation.  There is no need to
formalize a general probability inequality framework before these two files.

\paragraph{Cauchy--Schwarz.}
Use the weighted Lagrange identity from \Cref{lem:lagrange}.  The double sum of
squares is encoded by \code{MultiSums}; nonnegativity is a direct natural
statement.  This is more proof-friendly than a minimization or quadratic-
discriminant argument.

\paragraph{Independence and conditioning.}
For product laws, mixed-radix pairing and \code{MultiSums} prove factorization.
For conditioning, split on the natural cell mass.  A positive cell uses exact
cancellation; a zero cell uses \code{sum_zero_elim} to show all its weighted
terms vanish.  Sum the cell identities for total expectation.

\paragraph{Kolmogorov.}
Encode the first-crossing index by bounded search through the partial sums.
Prove disjointness of first-crossing indicators and factor the tail increment
using the product-law theorem.  The remaining argument is finite sum algebra.
This is the largest finite proof, so it should be attempted only after all
previous modules compile.

\subsection{Regular expectation API}

The file \code{lean/RegularExpectation.lean} defines nested rational intervals,
finite simple approximation levels, refinement maps preserving mass, an exact
refinement \(L^1\) distance, absolute-series names, and compact/improper
integrand interfaces.  The next implementation modules should be

\begin{verbatim}
IPRADerivations/Probability/RegularExpectation.lean
IPRADerivations/Probability/Series.lean
IPRADerivations/Probability/Integration.lean
IPRADerivations/Probability/Distributions.lean
IPRADerivations/Probability/Convergence.lean
IPRADerivations/Probability/Fourier.lean
\end{verbatim}

At this level the object-language sentence always contains a precision
variable.  A checker decodes a presentation and verifies one interval
inclusion with dyadic slack.  The mathematical proof computes the finite level
needed for that precision, invokes the finite kernel, and combines rational
error bounds.  It does not quantify over an ambient \code{Real}.

\subsection{Suggested milestone sequence}

\begin{enumerate}
\item Compile the ambient finite file against the repository's pinned Lean and
  Mathlib version, replacing every proof-package field by an actual theorem.
\item Add natural encoders and prove decoder/encoder round trips for valid
  codes.
\item Construct object-language derivations for normalization, constants,
  additivity, homogeneity, and monotonicity.
\item Add indicators, LOTUS, tail sums, and variance.
\item Add Markov, Chebyshev, and Cauchy--Schwarz.
\item Add product laws, conditioning, sample means, and Kolmogorov.
\item Only then implement regular names, series, quadrature, convergence, and
  Fourier modules.
\item Universally close each theorem and transport it with
  \code{pra_iff_native}; aggregate packages only after individual modules are
  checked.
\end{enumerate}

\subsection{Completion criterion}

The formalization is complete only when all of the following hold:
\begin{itemize}
\item the ambient files compile with the repository's pinned toolchain;
\item every decoder and primitive-recursive evaluator has proved defining
  equations;
\item each main theorem contains an explicit \code{Prf PRAAx} derivation;
\item no theorem is introduced through \code{axiom}, \code{sorry}, or a
  constant-zero checker masquerading as the semantic verifier;
\item the native proof is obtained through the existing proved bridge;
\item the Python validation remains outside the trusted proof boundary.
\end{itemize}
